% Compilar desde este directorio con:
% pdflatex Clase_02_Adquisicion_comprension_y_preparacion_inicial_de_datos.tex
\documentclass[aspectratio=169,11pt]{beamer}

\usepackage[utf8]{inputenc}
\usepackage[T1]{fontenc}
\usepackage[spanish,es-nodecimaldot]{babel}
\usepackage{amsmath}
\usepackage{amssymb}
\usepackage{booktabs}
\usepackage{tabularx}
\usepackage{array}
\usepackage{tikz}
\usetikzlibrary{arrows.meta,positioning}

\definecolor{UAPNavy}{HTML}{17324D}
\definecolor{UAPTeal}{HTML}{007F82}
\definecolor{UAPGold}{HTML}{E2A33A}
\definecolor{UAPLight}{HTML}{F3F6F8}
\definecolor{UAPGray}{HTML}{536471}

\tikzset{
  flow/.style={-{Stealth[length=2.4mm]},thick,draw=UAPNavy},
  datum/.style={draw=UAPGold,fill=UAPGold!16,rounded corners=2pt,align=center,minimum height=8mm},
  process/.style={draw=UAPTeal,fill=UAPTeal!12,rounded corners=3pt,very thick,align=center,minimum height=10mm},
  decision/.style={draw=UAPNavy,fill=UAPNavy!9,rounded corners=3pt,very thick,align=center,minimum height=10mm},
  card/.style={draw=UAPTeal,fill=UAPLight,rounded corners=3pt,align=center,minimum height=11mm}
}

\usetheme{Boadilla}
\usecolortheme{default}
\setbeamercolor{normal text}{fg=UAPNavy,bg=white}
\setbeamercolor{structure}{fg=UAPTeal}
\setbeamercolor{frametitle}{fg=white,bg=UAPNavy}
\setbeamercolor{title}{fg=white}
\setbeamercolor{subtitle}{fg=UAPGold}
\setbeamercolor{block title}{fg=white,bg=UAPTeal}
\setbeamercolor{block body}{fg=UAPNavy,bg=UAPLight}
\setbeamercolor{alerted text}{fg=UAPTeal}
\setbeamerfont{frametitle}{series=\bfseries}
\setbeamertemplate{navigation symbols}{}
\setbeamertemplate{itemize item}{\color{UAPGold}$\blacktriangleright$}
\setbeamertemplate{itemize subitem}{\color{UAPTeal}$\bullet$}
\setbeamertemplate{footline}{%
  \leavevmode%
  \hbox{%
    \begin{beamercolorbox}[wd=.78\paperwidth,ht=2.8ex,dp=1.2ex,leftskip=1em]{author in head/foot}%
      \color{UAPGray}\insertshorttitle
    \end{beamercolorbox}%
    \begin{beamercolorbox}[wd=.22\paperwidth,ht=2.8ex,dp=1.2ex,rightskip=1em plus1fil]{date in head/foot}%
      \color{UAPGray}\hfill Semana 2 \enspace | \enspace \insertframenumber/\inserttotalframenumber
    \end{beamercolorbox}%
  }%
}

\newcommand{\concepto}[1]{\textcolor{UAPTeal}{\textbf{#1}}}
\newcommand{\br}{\\}
\newcommand{\codigo}[1]{\texttt{\detokenize{#1}}}

\title[Data Science · Clase 2]{Adquisición, comprensión y preparación inicial de datos}
\subtitle{De fuentes heterogéneas a una tabla analítica defendible}
\author{Curso de Data Science}
\institute{Semana 2 · Clase 2}
\date{}

\begin{document}

% 1
\begin{frame}[plain]
  \begin{beamercolorbox}[wd=\paperwidth,ht=\paperheight,sep=2em]{frametitle}
    \vfill
    {\usebeamerfont{title}\usebeamercolor[fg]{title}\Huge\inserttitle\par}
    \vspace{0.5em}
    {\usebeamerfont{subtitle}\usebeamercolor[fg]{subtitle}\Large\insertsubtitle\par}
    \vspace{2em}
    {\color{white}\large\insertauthor\par}
    {\color{white}\insertinstitute\par}
    \vfill
  \end{beamercolorbox}
\end{frame}

% 2
\begin{frame}{Propósito y resultados de aprendizaje}
  \small
  \textbf{Propósito:} adquirir, comprender y preparar datos sin separar decisiones técnicas de su significado y uso.
  \begin{block}{Al finalizar, podremos}
    \begin{itemize}
      \item distinguir fuente, archivo, tabla, dataset y producto;
      \item documentar procedencia, unidad, granularidad y regla de una fila;
      \item evaluar calidad, faltantes, duplicados, claves y cardinalidad;
      \item integrar viajes, zonas y clima en una tabla \textbf{zona-hora};
      \item explicar pérdidas de agregación y límites del análisis.
    \end{itemize}
  \end{block}
  \centering Correspondencia: capítulos \textbf{2, 3 y 4} del libro.
\end{frame}

% 3
\begin{frame}{La preparación comienza con una pregunta}
  \begin{alertblock}{Punto de partida}
    ``Tenemos muchos archivos'' no define un problema de datos.
  \end{alertblock}
  Para describir entregas tardías por sucursal-día debemos precisar:
  \begin{itemize}
    \item producto esperado: tabla, gráfico o decisión;
    \item unidad final: sucursal-día;
    \item población, periodo y exclusiones;
    \item variables necesarias y momento de disponibilidad.
  \end{itemize}
  \centering\alert{La pregunta determina qué adquirir y cómo representarlo.}
\end{frame}

% 4
\begin{frame}{Fuente, archivo, tabla, dataset y producto}
  \small
  \begin{tabularx}{\textwidth}{>{\bfseries}p{.16\textwidth} p{.29\textwidth} X}
    \toprule
    Concepto & Función & Ejemplo breve \\
    \midrule
    Fuente & Sistema o institución de origen & portal municipal \\
    Archivo & Objeto de intercambio & \codigo{enero.parquet} \\
    Tabla & Estructura de filas y columnas & viajes \\
    Dataset & Colección delimitada y documentada & viajes + diccionario \\
    Producto & Evidencia destinada a un uso & perfil zona-hora \\
    \bottomrule
  \end{tabularx}
  \vfill
  \begin{block}{Distinción}
    Un dataset puede usar varios archivos; el producto es una transformación, no otro nombre para la fuente.
  \end{block}
\end{frame}

% 5
\begin{frame}{Procedencia y linaje}
  \centering
  \begin{tikzpicture}[node distance=4mm,scale=.82,transform shape]
    \node[decision] (f) {Fenómeno};
    \node[process,right=of f] (c) {Captura};
    \node[datum,right=of c] (s) {Fuente};
    \node[process,right=of s] (e) {Extracción};
    \node[process,right=of e] (t) {Transformación};
    \node[datum,right=of t] (p) {Producto};
    \foreach \x/\y in {f/c,c/s,s/e,e/t,t/p}{\draw[flow] (\x)--(\y);}
  \end{tikzpicture}
  \vspace{.6em}
  \begin{itemize}
    \item Registrar responsable, propósito, URL, acceso, versión y licencia.
    \item Conservar filtros, conversiones, uniones y artefactos resultantes.
  \end{itemize}
  \begin{block}{Dos preguntas}
    Procedencia: ¿de dónde viene? \quad Linaje: ¿qué le hicimos?
  \end{block}
\end{frame}

% 6
\begin{frame}{Unidad y granularidad}
  \begin{columns}[T]
    \column{.48\textwidth}
      \begin{block}{Unidad de análisis}
        Entidad elemental sobre la que se describe o compara.
      \end{block}
    \column{.48\textwidth}
      \begin{block}{Granularidad}
        Nivel de detalle temporal, espacial y de entidad.
      \end{block}
  \end{columns}
  \vspace{.5em}
  \small
  \begin{tabularx}{\textwidth}{>{\bfseries}p{.24\textwidth} X X}
    \toprule
    Representación & Unidad & Granularidad \\
    \midrule
    ventas & transacción & producto-instante \\
    resumen & tienda-día & tienda $\times$ día \\
    \bottomrule
  \end{tabularx}
  \vfill
  \centering\alert{Cambiar granularidad cambia preguntas válidas, dependencias y filas.}
\end{frame}

% 7
\begin{frame}{La regla de una fila}
  \begin{block}{Completar antes de limpiar}
    \centering\large Una fila representa \rule{5cm}{.4pt}.
  \end{block}
  Ejemplo: ``una medición de un sensor en un instante'', no ``un sensor''.
  \begin{itemize}
    \item ¿Qué columnas identifican esa unidad?
    \item ¿Una clave repetida es repetición, corrección o duplicado?
    \item ¿Listas, intervalos o medidas repetidas rompen la regla?
  \end{itemize}
  \begin{alertblock}{Riesgo}
    Sin esta regla, \codigo{drop_duplicates()} puede borrar observaciones legítimas.
  \end{alertblock}
\end{frame}

% 8
\begin{frame}{Tipos físicos y tipos semánticos}
  \small
  \begin{tabularx}{\textwidth}{>{\ttfamily}p{.19\textwidth} p{.22\textwidth} X}
    \toprule
    \normalfont\textbf{Valor} & \textbf{Tipo físico} & \textbf{Tipo semántico correcto} \\
    \midrule
    00123 & texto & identificador nominal \\
    2024-01-05 & texto & fecha \\
    1, 2, 3 & entero & categoría ordinal \\
    -73.98 & decimal & longitud geográfica \\
    \bottomrule
  \end{tabularx}
  \vfill
  \begin{block}{Principio}
    El tipo físico permite operaciones; el semántico determina cuáles tienen sentido.
  \end{block}
  \centering Promediar códigos postales es computable, pero inválido.
\end{frame}

% 9
\begin{frame}{Calidad apta al propósito}
  \small
  \begin{tabularx}{\textwidth}{>{\bfseries}p{.22\textwidth} X}
    \toprule
    Dimensión & Pregunta de control \\
    \midrule
    Completitud & ¿Faltan campos o unidades necesarias? \\
    Validez & ¿Cumple dominio, formato y rango? \\
    Consistencia & ¿Concuerdan campos y fuentes? \\
    Unicidad & ¿La clave identifica una fila? \\
    Oportunidad & ¿Llega antes de usarla? \\
    Cobertura & ¿Representa población, lugar y periodo? \\
    \bottomrule
  \end{tabularx}
  \vfill
  \centering\alert{Un dato puede ser correcto y aun así no ser apto.}
\end{frame}

% 10
\begin{frame}{Faltantes: ausencia con significado}
  Un nulo puede significar:
  \begin{itemize}
    \item no medido, no aplicable o desconocido;
    \item fallo de captura o unión sin correspondencia;
    \item valor censurado o todavía no disponible.
  \end{itemize}
  \begin{block}{Diagnóstico mínimo}
    Proporción por variable, patrón por grupo y periodo, y relación con el proceso de captura.
  \end{block}
  \centering\alert{No imputar hasta explicar el mecanismo y el costo de equivocarse.}
\end{frame}

% 11
\begin{frame}{Duplicados: definir antes de eliminar}
  \begin{columns}[T]
    \column{.32\textwidth}\begin{block}{Exacto}Todas las columnas coinciden.\end{block}
    \column{.32\textwidth}\begin{block}{De clave}Coincide la clave; difieren atributos.\end{block}
    \column{.32\textwidth}\begin{block}{De entidad}Mismo objeto sin clave común.\end{block}
  \end{columns}
  \vspace{.6em}
  Ejemplo: dos lecturas del mismo paciente-hora pueden ser duplicados, reintentos o mediciones válidas.
  \begin{alertblock}{Decisión documentada}
    Conservar, consolidar, versionar o excluir según la regla de una fila y el proceso generador.
  \end{alertblock}
\end{frame}

% 12
\begin{frame}{Claves y cardinalidad}
  \begin{itemize}
    \item \concepto{Clave primaria}: identifica filas en una tabla.
    \item \concepto{Clave foránea}: referencia otra tabla.
    \item \concepto{Cardinalidad}: \codigo{one_to_one}, \codigo{many_to_one}, \codigo{one_to_many} o \codigo{many_to_many}.
  \end{itemize}
  \begin{block}{Ejemplo}
    Muchas ventas pertenecen a una tienda: ventas $\rightarrow$ tiendas es \textbf{many-to-one}.
  \end{block}
  \centering\alert{Una unión many-to-many accidental multiplica filas y totales.}
\end{frame}

% 13
\begin{frame}{Integración segura}
  \begin{enumerate}
    \item Declarar unidad y clave de cada tabla.
    \item Verificar unicidad del lado que debe ser ``uno''.
    \item Validar cardinalidad en el código.
    \item Medir coincidencias y no coincidencias.
    \item Comprobar filas, unidades y totales antes y después.
  \end{enumerate}
  \vfill
  \begin{block}{Principio}
    \centering Toda unión es una hipótesis sobre correspondencia.
  \end{block}
\end{frame}

% 14
\begin{frame}{Integración temporal}
  \centering
  \begin{tikzpicture}[node distance=8mm,scale=.9,transform shape]
    \node[datum] (ev) {Tiempo del evento};
    \node[process,right=of ev] (reg) {Tiempo de registro};
    \node[decision,right=of reg] (disp) {Disponibilidad};
    \draw[flow] (ev)--(reg); \draw[flow] (reg)--(disp);
  \end{tikzpicture}
  \vspace{.7em}
  \begin{itemize}
    \item Interpretar zona horaria de origen y convertir a una común.
    \item Definir intervalos, bordes y horario de verano.
    \item Agregar fuentes a granularidad compatible.
  \end{itemize}
  \begin{alertblock}{No equivalencia}
    \centering 10:00 UTC $\neq$ 10:00 America/New\_York
  \end{alertblock}
\end{frame}

% 15
\begin{frame}{Integración espacial}
  Una relación espacial puede usar:
  \begin{itemize}
    \item identificador administrativo estable;
    \item punto dentro de polígono;
    \item zona más cercana o intersección.
  \end{itemize}
  \begin{block}{Registrar siempre}
    Sistema de referencia de coordenadas (CRS), versión de límites y casos sin zona.
  \end{block}
  \centering\alert{Un centroide resume geometría; no reemplaza el polígono.}
\end{frame}

% 16
\begin{frame}{Agregación y pérdida de información}
  \[
    n_g=\sum_i \mathbb{1}(G_i=g), \qquad
    \bar{x}_g=\frac{1}{n_g}\sum_{i:G_i=g}x_i
  \]
  \begin{columns}[T]
    \column{.48\textwidth}
      \begin{block}{Ganamos}
        Comparabilidad, menor volumen y unidad útil para la pregunta.
      \end{block}
    \column{.48\textwidth}
      \begin{block}{Perdemos}
        Orden intra-grupo, detalle individual y reconstrucción de casos.
      \end{block}
  \end{columns}
  \vfill
  \centering Conservar conteo, dispersión, cobertura y regla junto a la media.
\end{frame}

% 17
\begin{frame}{Reproducibilidad desde la adquisición}
  Una ejecución reconstruible registra:
  \begin{itemize}
    \item URL, fecha de acceso, periodo y parámetros;
    \item versiones de librerías y formato;
    \item código ordenado de inicio a fin;
    \item validaciones y decisiones de exclusión;
    \item salida con esquema, clave y procedencia.
  \end{itemize}
  \begin{block}{Notebook de la clase}
    \centering\codigo{Taller_Clase_02_Movilidad.ipynb}
  \end{block}
\end{frame}

% 18
\begin{frame}{Del concepto al taller}
  \begin{block}{Pregunta descriptiva}
    ¿Cómo varían los viajes Yellow Taxi reportados por zona de origen y hora, y qué asociación muestran con el clima?
  \end{block}
  \centering
  \begin{tikzpicture}[node distance=6mm,scale=.86,transform shape]
    \node[datum] (v) {TLC viajes};
    \node[datum,below=of v] (z) {Zonas + shapefile};
    \node[process,right=of v,yshift=-5mm] (a) {Zona-hora};
    \node[datum,right=of a] (c) {NOAA clima};
    \draw[flow] (v)--(a); \draw[flow] (z)--(a); \draw[flow] (c)--(a);
  \end{tikzpicture}
  \vfill
  \alert{Viajes reportados por TLC $\neq$ demanda total de movilidad.}
\end{frame}

% 19
\begin{frame}{Taller 1: fuentes y contrato analítico}
  \scriptsize
  \textbf{Notebook:} \codigo{Taller_Clase_02_Movilidad.ipynb}
  \begin{tabularx}{\textwidth}{>{\bfseries}p{.15\textwidth} p{.43\textwidth} X}
    \toprule
    Fuente & Archivo / tabla & Unidad original \\
    \midrule
    NYC TLC & Yellow Taxi enero 2024, Parquet & viaje reportado \\
    NYC TLC & \codigo{taxi_zone_lookup.csv} & zona-ID \\
    NYC TLC & \codigo{taxi_zones.zip}, shapefile & polígono-zona \\
    NOAA GHCNh & estación USW00094728, PSV & estación-observación \\
    \bottomrule
  \end{tabularx}
  \vfill
  \begin{block}{Producto}
    Una fila = \textbf{zona de origen-hora local}; clave: \codigo{PULocationID, hora}.
  \end{block}
\end{frame}

% 20
\begin{frame}[fragile]{Taller 2: construir la URL TLC}
\begin{verbatim}
def construir_url_tlc(tipo, anio, mes):
    prefijos = {
        "yellow": "yellow_tripdata",
        "green": "green_tripdata",
        "fhv": "fhv_tripdata",
        "fhvhv": "fhvhv_tripdata",
    }
    return ("https://d37ci6vzurychx.cloudfront.net/trip-data/"
            f"{prefijos[tipo]}_{anio}-{mes:02d}.parquet")
\end{verbatim}
  \small Parámetros explícitos hacen visible el periodo y evitan editar URL manualmente.
\end{frame}

% 21
\begin{frame}[fragile]{Taller 3: cargar Parquet con proyección}
\begin{verbatim}
cols = ["tpep_pickup_datetime", "tpep_dropoff_datetime",
        "PULocationID", "DOLocationID", "passenger_count",
        "trip_distance", "fare_amount", "total_amount"]
url = construir_url_tlc("yellow", 2024, 1)
viajes = pd.read_parquet(url, columns=cols)
print(viajes.shape)
\end{verbatim}
  \begin{block}{Procedencia}
    Fuente pública y secundaria, generada con propósito operativo/administrativo.
  \end{block}
\end{frame}

% 22
\begin{frame}[fragile]{Taller 4: inspeccionar antes de transformar}
\begin{verbatim}
viajes.info()
display(viajes.head(3))
display(viajes.describe(include="all").T)
print(viajes.dtypes)
\end{verbatim}
  \begin{itemize}
    \item ¿Qué representa una fila y qué rango temporal aparece?
    \item ¿Los ID son medidas o categorías?
    \item ¿Hay importes o duraciones imposibles?
  \end{itemize}
  \centering\alert{La inspección genera hipótesis; los controles las verifican.}
\end{frame}

% 23
\begin{frame}[fragile]{Taller 5: controles de calidad básicos}
\begin{verbatim}
clave = ["tpep_pickup_datetime", "tpep_dropoff_datetime",
         "PULocationID", "DOLocationID", "total_amount"]
qc = {
    "nulos_pct": viajes.isna().mean().mul(100),
    "duplicados_clave": viajes.duplicated(clave).sum(),
    "duracion_no_positiva": (viajes.tpep_dropoff_datetime <=
                              viajes.tpep_pickup_datetime).sum(),
}
\end{verbatim}
  \small La clave es una aproximación operativa, no un ID oficial de viaje.
\end{frame}

% 24
\begin{frame}[fragile]{Taller 6: validez y filtros explícitos}
\begin{verbatim}
validos = (
    viajes["PULocationID"].notna()
    & (viajes["trip_distance"] >= 0)
    & (viajes["total_amount"] >= 0)
    & (viajes["tpep_dropoff_datetime"] >
       viajes["tpep_pickup_datetime"])
)
viajes_limpios = viajes.loc[validos].copy()
\end{verbatim}
  \small El catálogo TLC validará los ID al unir. Plausibilidad no equivale a verdad.
\end{frame}

% 25
\begin{frame}[fragile]{Taller 7: zonas y unión many-to-one}
\begin{verbatim}
zones_url = ("https://d37ci6vzurychx.cloudfront.net/"
             "misc/taxi_zone_lookup.csv")
zonas = pd.read_csv(zones_url)
assert zonas["LocationID"].is_unique
viajes_z = viajes_limpios.merge(
    zonas, left_on="PULocationID", right_on="LocationID",
    how="left", validate="many_to_one", indicator=True)
\end{verbatim}
  \small Auditar \codigo{_merge}: cada viaje encuentra como máximo una zona.
\end{frame}

% 26
\begin{frame}[fragile]{Taller 8: cobertura y pérdidas de la unión}
\begin{verbatim}
cobertura = viajes_z["_merge"].value_counts(dropna=False)
sin_zona = viajes_z.loc[
    viajes_z["_merge"] == "left_only", "PULocationID"
].value_counts()
assert len(viajes_z) == len(viajes_limpios)
print(cobertura, sin_zona.head())
\end{verbatim}
  \begin{block}{Interpretación}
    Una unión izquierda conserva viajes, pero puede dejar atributos nulos.
  \end{block}
\end{frame}

% 27
\begin{frame}[fragile]{Taller 9: shapefile y geometría}
\begin{verbatim}
geo_url = ("https://d37ci6vzurychx.cloudfront.net/"
           "misc/taxi_zones.zip")
zonas_geo = gpd.read_file(f"zip://{geo_url}")
zonas_geo["LocationID"] = zonas_geo["LocationID"].astype(int)
zonas_geo = zonas_geo.to_crs(4326)
assert zonas_geo["LocationID"].is_unique
\end{verbatim}
  \small El lookup aporta nombres; el shapefile, polígonos. Registrar CRS y versión.
\end{frame}

% 28
\begin{frame}[fragile]{Taller 10: construir zona-hora}
\begin{verbatim}
viajes_z["hora"] = (
    viajes_z["tpep_pickup_datetime"]
    .dt.tz_localize("America/New_York", ambiguous="NaT",
                    nonexistent="shift_forward")
    .dt.floor("h")
)
zona_hora = (viajes_z.groupby(["PULocationID", "hora"])
  .agg(viajes_reportados=("PULocationID", "size"),
       distancia_media=("trip_distance", "mean")).reset_index())
\end{verbatim}
  \small Cambia la unidad de viaje a zona-hora y elimina detalle individual.
\end{frame}

% 29
\begin{frame}[fragile]{Taller 11: NOAA e interpretación UTC}
\begin{verbatim}
meteo_url = ("https://www.ncei.noaa.gov/oa/global-historical-"
  "climatology-network/hourly/access/by-year/2024/psv/"
  "GHCNh_USW00094728_2024.psv")
clima = pd.read_csv(meteo_url, sep="|", low_memory=False)
clima["DATE_UTC"] = pd.to_datetime(clima["DATE"], utc=True)
clima["DATE_NY"] = clima["DATE_UTC"].dt.tz_convert(
    "America/New_York")
\end{verbatim}
  \small Convertimos un instante UTC; no solo la etiqueta de una hora ingenua.
\end{frame}

% 30
\begin{frame}[fragile]{Taller 12: clima horario de enero}
\begin{verbatim}
inicio = pd.Timestamp("2024-01-01", tz="America/New_York")
fin = pd.Timestamp("2024-02-01", tz="America/New_York")
clima_mes = clima[clima["DATE_NY"].between(
    inicio, fin, inclusive="left")].copy()
clima_mes["hora"] = clima_mes["DATE_NY"].dt.floor("h")
vars_num = ["temperature", "relative_humidity",
            "wind_speed", "precipitation"]
clima_h = clima_mes.groupby("hora")[vars_num].mean().reset_index()
\end{verbatim}
  \small Promediar muestras subhorarias pierde extremos y secuencia.
\end{frame}

% 31
\begin{frame}[fragile]{Taller 13: unir zona-hora con clima}
\begin{verbatim}
assert clima_h["hora"].is_unique
analitica = zona_hora.merge(
    clima_h, on="hora", how="left",
    validate="many_to_one", indicator="union_clima")
print(analitica["union_clima"].value_counts())
print(analitica.isna().mean()
      .sort_values(ascending=False).head())
\end{verbatim}
  \small Una estación para muchas zonas supone representatividad espacial.
\end{frame}

% 32
\begin{frame}[fragile]{Taller 14: perfil horario de viajes reportados}
\begin{verbatim}
perfil_hora = (analitica.assign(
    hora_dia=analitica["hora"].dt.hour)
    .groupby("hora_dia", as_index=False)
    ["viajes_reportados"].sum())
plt.plot(perfil_hora["hora_dia"],
         perfil_hora["viajes_reportados"], marker="o")
plt.ylabel("Viajes Yellow Taxi reportados")
plt.xlabel("Hora local")
\end{verbatim}
  \small Describe registros TLC de enero de 2024, no demanda total.
\end{frame}

% 33
\begin{frame}[fragile]{Taller 15: mapa por zona}
\begin{verbatim}
por_zona = (analitica.groupby("PULocationID", as_index=False)
            ["viajes_reportados"].sum())
mapa = zonas_geo.merge(
    por_zona, left_on="LocationID", right_on="PULocationID",
    how="left", validate="one_to_one")
mapa.plot(column="viajes_reportados", cmap="YlGnBu",
          legend=True, missing_kwds={"color": "lightgrey"})
\end{verbatim}
  \small Los conteos mezclan actividad, tamaño, cobertura y registro.
\end{frame}

% 34
\begin{frame}[fragile]{Taller 16: clima y asociación descriptiva}
\begin{verbatim}
por_hora = (analitica.groupby("hora", as_index=False)
  .agg(viajes_reportados=("viajes_reportados", "sum"),
       temperatura=("temperature", "first"),
       precipitacion=("precipitation", "first")))
plt.scatter(por_hora["temperatura"], por_hora["viajes_reportados"],
            c=por_hora["precipitacion"], cmap="Blues")
r = por_hora[
    ["temperatura", "viajes_reportados"]
].corr().iloc[0, 1]
\end{verbatim}
  \small Hora y calendario pueden confundir. \alert{Correlación no implica causalidad.}
\end{frame}

% 35
\begin{frame}[fragile]{Taller 17: guardar un producto reproducible}
\begin{verbatim}
analitica = analitica.sort_values(["hora", "PULocationID"])
assert not analitica.duplicated(
    ["PULocationID", "hora"]).any()
alcance = "semana_1_enero_2024" if MODO_CLASE else "enero_2024"
ruta = Path(tempfile.gettempdir()) / f"zona_hora_{alcance}.parquet"
analitica.to_parquet(ruta, index=False)
metadatos = {"unidad": "zona de origen-hora local",
             "zona_horaria": "America/New_York",
             "fuentes": [url, zones_url, geo_url, meteo_url]}
\end{verbatim}
  \small Entregar controles, exclusiones, versiones y fecha de ejecución.
\end{frame}

% 36
\begin{frame}{Producto y evidencia de aprendizaje}
  \small
  \textbf{Entregable:} \codigo{Taller_Clase_02_Movilidad.ipynb} ejecutado de inicio a fin.
  \begin{itemize}
    \item inventario de cuatro fuentes y procedencia;
    \item regla de una fila y diccionario semántico mínimo;
    \item reporte de faltantes, duplicados, validez, cobertura y cardinalidad;
    \item tabla zona-hora temporal con alcance explícito y clave única;
    \item perfil horario, mapa y gráfico clima-viajes;
    \item cinco hallazgos con evidencia y tres límites.
  \end{itemize}
  \begin{block}{Criterio}
    Otra persona puede ejecutar, auditar y explicar cada transformación.
  \end{block}
\end{frame}

% 37
\begin{frame}{Límites y afirmaciones permitidas}
  Podemos describir \textbf{viajes Yellow Taxi reportados} y asociaciones en el periodo y cobertura analizados.
  \vspace{.4em}
  \textbf{No podemos afirmar directamente:}
  \begin{itemize}
    \item demanda total de movilidad o solicitudes no atendidas;
    \item comportamiento de todas las plataformas, personas o meses;
    \item causalidad del clima sobre los viajes;
    \item clima idéntico en toda zona por usar una estación;
    \item exactitud de cada registro por superar controles de plausibilidad.
  \end{itemize}
\end{frame}

% 38
\begin{frame}{Cierre: de datos disponibles a evidencia defendible}
  \small
  \begin{enumerate}
    \item Formular propósito y producto antes de descargar.
    \item Declarar fuente, procedencia, unidad y regla de una fila.
    \item Tratar tipos, faltantes y duplicados según significado.
    \item Validar claves, cardinalidad y cobertura al integrar.
    \item Alinear tiempo y espacio antes de agregar.
    \item Documentar pérdidas, límites y cadena reproducible.
  \end{enumerate}
  \vfill
  \begin{block}{Para continuar}
    Capítulos 2, 3 y 4. Próximo paso: exploración estadística sin exceder lo que los datos representan.
  \end{block}
\end{frame}

\end{document}
